\chapter{Observability}

\label{ch:observability}

An operator asks three questions of a running pii, and each has one surface that answers it
faster than the others: \emph{is it up and is it the build I think it is}, \emph{what did this
run actually do}, and \emph{why was that refused}. This chapter covers those surfaces in that
order --- the daemon's account of itself, the event stream, the metric set, the run journal,
and the audit trail --- and ends with the correlation order that turns five commands into one
answer.

The figures and tables are generated from the code: the event names from
\texttt{crates/pii-observe/src/names.rs}, the metric names from the same file, the journal
shape from \texttt{crates/pii-abi/src/journal.rs}, and the command output from
\texttt{crates/pii-cli/src/cmd\_health.rs} and \texttt{cmd\_metrics.rs}.

\section{Four surfaces, four questions}

\begin{center}
\begin{minipage}{\linewidth}
\centering
\resizebox{\linewidth}{!}{%
\begin{tikzpicture}[node distance=8mm and 12mm]
  \node[v] (q) {what do I need?};
  \node[b, above right=8mm and 14mm of q] (up) {is it up?};
  \node[b, right=14mm of q, yshift=2mm] (did) {what happened?};
  \node[b, below right=8mm and 14mm of q] (no) {why refused?};
  \node[b, right=of up] (health) {\texttt{pii health}\\\scriptsize identity, stores, mounts};
  \node[b, right=of did] (trace) {\texttt{pii run trace}\\\scriptsize the journal};
  \node[b, right=of no] (audit) {\texttt{audit.log}\\\scriptsize denials only};
  \node[b, below=of did] (loop) {\texttt{pii loop events}\\\scriptsize the autonomy loop};
  \draw[f] (q) -- (up);
  \draw[f] (q) -- (did);
  \draw[f] (q) -- (no);
  \draw[f] (up) -- (health);
  \draw[f] (did) -- (trace);
  \draw[f] (no) -- (audit);
  \draw[sf] (did) -- (loop);
\end{tikzpicture}}
\captionof{figure}{The entry point for each question. \texttt{pii health} is the right first command for all three, because it tells you whether the other surfaces can answer at all.}
\end{minipage}
\end{center}

\begin{center}
\begin{tabular}{@{}>{\raggedright\arraybackslash}p{0.20\textwidth}>{\raggedright\arraybackslash}p{0.28\textwidth}>{\raggedright\arraybackslash}p{0.42\textwidth}@{}}
  \toprule
  \textbf{Surface} & \textbf{Written by} & \textbf{Good for, and its limit} \\
  \midrule
  \texttt{pii health} & the daemon, on request & identity, config digest, stores, mounts,
    counters since deploy. It is a snapshot, not a history. \\
  the event log & the daemon, per event & everything the daemon decided. Off unless
    \texttt{-{}-log-file} or \texttt{[core].log\_file} is set. \\
  \texttt{audit.log} & the daemon, per denial & the security-relevant subset. Always on unless
    \texttt{-{}-audit off}. Never contains a kernel-level denial raised inside a child. \\
  the run journal & the engine, per run & what a run did, node by node, replayable. Grows with
    runs; retention applies. \\
  \bottomrule
\end{tabular}
\end{center}

\section{The anatomy of one event}

Every event is one JSON object on one line. The file sink adds a wall-clock stamp; the rest
of the object is the event itself.

\begin{center}
\begin{minipage}{\linewidth}
\centering
\resizebox{\linewidth}{!}{%
\begin{tikzpicture}[node distance=8mm and 12mm]
  \node[b] (ts) {\texttt{ts\_ms}\\\scriptsize wall clock, added by the sink};
  \node[b, right=of ts] (sub) {\texttt{subsystem}\\\scriptsize who reported};
  \node[b, right=of sub] (lvl) {\texttt{level}\\\scriptsize trace..error};
  \node[b, right=of lvl] (kind) {\texttt{kind}\\\scriptsize the name};
  \node[b, below=of kind] (fields) {\texttt{fields}\\\scriptsize structured detail};
  \node[b, below=of lvl] (span) {\texttt{span}\\\scriptsize the step};
  \node[b, below=of sub] (run) {\texttt{run}, \texttt{session}\\\scriptsize optional};
  \draw[sf] (ts) -- (sub);
  \draw[sf] (sub) -- (run);
  \draw[sf] (lvl) -- (span);
  \draw[sf] (kind) -- (fields);
\end{tikzpicture}}
\captionof{figure}{Six fields are always present (\texttt{ts\_ms}, \texttt{subsystem}, \texttt{level}, \texttt{kind}, \texttt{fields}, \texttt{span}); \texttt{run} and \texttt{session} appear only when the event belongs to one.}
\end{minipage}
\end{center}

The serialised form, with the run and session fields omitted because the event has neither:

\begin{lstlisting}[language=bash]
{"subsystem":"x","level":"info","kind":"k","fields":{"a":1},"span":null}
\end{lstlisting}

\texttt{level} is one of \texttt{trace}, \texttt{debug}, \texttt{info}, \texttt{warn},
\texttt{error}. In practice the daemon emits \texttt{info} for lifecycle, \texttt{warn} for
a refusal and \texttt{error} for a failure it could not prevent; \texttt{trace} and
\texttt{debug} exist in the type and are not the ones to grep for during an incident.

The two correlation fields are the reason a line can be trusted. \texttt{run} is the run id
and \texttt{session} the session id, so one \texttt{grep} answers \enquote{everything that
happened to this run} without matching a node name that three graphs share.

\section{The name registry}

Names are not string literals at call sites. Every event and every metric is a constant in one
file, and the file also holds an \texttt{ALL} array that a unit test checks for duplicates and
empty names. The count is exact and worth knowing, because it is the size of the vocabulary an
operator has to recognise:

\begin{center}
\begin{tabular}{@{}>{\raggedright\arraybackslash}p{0.30\textwidth}>{\raggedright\arraybackslash}p{0.16\textwidth}>{\raggedright\arraybackslash}p{0.44\textwidth}@{}}
  \toprule
  \textbf{Kind of name} & \textbf{Count} & \textbf{Where it is used} \\
  \midrule
  event names & 150 & the \texttt{kind} field of a log line, and the audit filter \\
  metric names & 40 & the series names \texttt{pii metrics} prints \\
  the \texttt{ALL} registry & 1 & the duplicate check, not emitted \\
  \midrule
  constants in the file & 191 & 150 + 40 + 1 \\
  \bottomrule
\end{tabular}
\end{center}

The 150 events fall into 38 subsystems. The largest are the ones an installation touches
every day: channels, learning, models, memory, runs.

\begin{center}
\begin{tabular}{@{}>{\raggedright\arraybackslash}p{0.16\textwidth}r@{\hspace{6mm}}>{\raggedright\arraybackslash}p{0.16\textwidth}r@{\hspace{6mm}}>{\raggedright\arraybackslash}p{0.16\textwidth}r@{}}
  \toprule
  \textbf{Subsystem} & \textbf{n} & \textbf{Subsystem} & \textbf{n} & \textbf{Subsystem} & \textbf{n} \\
  \midrule
  \texttt{channel} & 14 & \texttt{model} & 11 & \texttt{daemon} & 7 \\
  \texttt{learn} & 13 & \texttt{memory} & 10 & \texttt{mqtt} & 6 \\
  \texttt{run} & 10 & \texttt{coder} & 9 & \texttt{hal} & 6 \\
  \texttt{wasm} & 8 & \texttt{worker} & 5 & \texttt{trajectory} & 4 \\
  \texttt{proto} & 4 & \texttt{sop} & 4 & \texttt{graph} & 3 \\
  \texttt{net} & 3 & \texttt{cron} & 3 & \texttt{evolve} & 3 \\
  \texttt{timer} & 3 & \texttt{tools} & 2 & \texttt{queue} & 2 \\
  \texttt{tui} & 2 & \texttt{web} & 2 & \texttt{fetch} & 2 \\
  \texttt{effect} & 1 & \texttt{store} & 1 & \texttt{core} & 1 \\
  \texttt{health} & 1 & \texttt{observe} & 1 & \texttt{criteria} & 1 \\
  \texttt{skill} & 1 & \texttt{cli} & 1 & \texttt{lint} & 1 \\
  \texttt{cursor} & 1 & \texttt{mcp} & 1 & \texttt{node} & 1 \\
  \texttt{log} & 1 & \texttt{uds} & 1 & & \\
  \bottomrule
\end{tabular}
\end{center}

A few of the names are load-bearing enough to memorise, because they name the failures that
are otherwise silent:

\begin{center}\footnotesize
\begin{tabular}{@{}>{\raggedright\arraybackslash}p{0.34\textwidth}>{\raggedright\arraybackslash}p{0.56\textwidth}@{}}
  \toprule
  \textbf{Name} & \textbf{What it means when you see it} \\
  \midrule
  \texttt{graph.node.dead} & a node was marked dead; the run cannot proceed through it \\
  \texttt{graph.join.unsatisfiable} & a join can never be satisfied; the graph is wrong, not
    the run \\
  \texttt{observe.events.dropped} & the event pipeline dropped events; your log has a gap \\
  \texttt{core.cap.denied} & a capability the node did not hold was requested \\
  \texttt{store.persist.error} & a write to a store failed; the run's account is incomplete \\
  \texttt{model.provider.auth\_failed} & the key, not the model, is the problem \\
  \texttt{model.provider.rate\_limited} & the provider refused for quota; retrying now is waste \\
  \texttt{coder.env.refused} & a variable named in \texttt{env\_pass} was not passed \\
  \bottomrule
\end{tabular}
\end{center}

\section{Metrics}

\texttt{pii metrics} asks the daemon for its status payload and renders gauges and counters in
the text exposition format, one series per line, with labels in braces:

\begin{lstlisting}[language=bash]
pii_store_bytes{store="main"} 1048576
pii_runs 42
pii_run_latency_ms{outcome="succeeded",quantile="p95"} 132
\end{lstlisting}

It needs a running daemon: with none, it prints \enquote{no daemon answered} and exits 2. The
same numbers are in \texttt{pii health -{}-json}, which is the better choice when you want to
keep the payload for a record rather than read it.

\begin{center}
\begin{minipage}{\linewidth}
\centering
\resizebox{\linewidth}{!}{%
\begin{tikzpicture}[node distance=8mm and 12mm]
  \node[b] (sys) {\texttt{sys.status}\\\scriptsize one payload};
  \node[b, above right=6mm and 12mm of sys] (store) {store family\\\scriptsize bytes, rows, growth};
  \node[b, right=12mm of sys] (runs) {run family\\\scriptsize counts, latency, uptime};
  \node[b, below right=6mm and 12mm of sys] (sub) {subsystem family\\\scriptsize mqtt, channel, hal, wasm};
  \node[b, right=12mm of runs] (learn) {learning family\\\scriptsize candidates, verdicts, promotions};
  \node[b, right=12mm of sub] (mem) {memory family\\\scriptsize docs, entities, tiers, archives};
  \draw[f] (sys) -- (store);
  \draw[f] (sys) -- (runs);
  \draw[f] (sys) -- (sub);
  \draw[sf] (runs) -- (learn);
  \draw[sf] (sub) -- (mem);
\end{tikzpicture}}
\captionof{figure}{Six families, one payload. Every metric is derived from state the daemon already holds, so asking for metrics never changes what it is doing.}
\end{minipage}
\end{center}

The full set, with the label each series carries. \emph{Gauge} means a value that can go down;
\emph{counter} means one that only grows until the daemon restarts.

{\footnotesize
\begin{longtable}{@{}>{\raggedright\arraybackslash}p{0.33\textwidth}>{\raggedright\arraybackslash}p{0.08\textwidth}>{\raggedright\arraybackslash}p{0.47\textwidth}@{}}
  \toprule
  \textbf{Series} & \textbf{Type} & \textbf{Meaning} \\
  \midrule
  \endhead
  \texttt{pii\_runs} & gauge & runs in the main store \\
  \texttt{pii\_sessions} & gauge & sessions in the main store \\
  \texttt{pii\_approvals\_pending} & gauge & approvals awaiting a decision \\
  \texttt{pii\_journal\_events\_total} & counter & journal events ever appended \\
  \texttt{pii\_memory\_docs} & gauge & memory documents \\
  \texttt{pii\_cron\_entries} & gauge & configured cron entries \\
  \texttt{pii\_store\_bytes} & gauge, \texttt{store} & size of a store file in bytes \\
  \texttt{pii\_store\_format\_version} & gauge, \texttt{store} & persisted format version \\
  \texttt{pii\_rss\_bytes\_per\_run} & gauge & RSS growth per run over the recent window \\
  \texttt{pii\_process\_rss\_kb} & gauge & resident set size in KiB \\
  \texttt{pii\_store\_bytes\_per\_run} & gauge & main-store bytes divided by runs \\
  \texttt{pii\_run\_latency\_ms} & gauge, \texttt{outcome}, \texttt{quantile} & run latency,
    upper bucket bound \\
  \texttt{pii\_run\_latency\_count} & counter, \texttt{outcome} & runs in the histogram \\
  \texttt{pii\_uptime\_ms} & gauge & daemon uptime \\
  \texttt{pii\_mqtt\_received\_total} & counter & messages accepted from brokers \\
  \texttt{pii\_mqtt\_dropped\_total} & counter & messages dropped: queue full, not UTF-8 \\
  \texttt{pii\_mqtt\_published\_total} & counter & messages published \\
  \texttt{pii\_worker\_requests\_refused\_total} & counter & refusals around a drain or reload \\
  \texttt{pii\_channel\_posts\_total} & counter & messages a \texttt{post} node delivered \\
  \texttt{pii\_channel\_post\_refusals\_total} & counter & \texttt{post} calls refused or failed \\
  \texttt{pii\_quarantine\_refusals\_total} & counter & runs refused by a promotion quarantine \\
  \texttt{pii\_quarantine\_refusals\_recorded\_total} & counter & refusals whose deny reached the
    ledger; a gap to the line above is a finding \\
  \texttt{pii\_hal\_ops\_total} & counter, \texttt{kind} & HAL operations that succeeded \\
  \texttt{pii\_hal\_refusals\_total} & counter & HAL operations refused by claim or scope \\
  \texttt{pii\_wasm\_calls\_total} & counter, \texttt{outcome} & wasm calls by outcome \\
  \texttt{pii\_wasm\_refusals\_total} & counter & modules or calls refused before running \\
  \texttt{pii\_learn\_candidates} & gauge, \texttt{kind} & candidates from the last scan \\
  \texttt{pii\_learn\_refusals\_total} & counter & candidates refused for missing evidence \\
  \texttt{pii\_learn\_candidates\_proposed\_total} & counter & distinct candidates the unattended
    loop proposed \\
  \texttt{pii\_learn\_unattended\_runs\_total} & counter & unattended extract+verify passes \\
  \texttt{pii\_learn\_verdicts} & gauge, \texttt{verdict} & verdicts of the last verification \\
  \texttt{pii\_learn\_promotions} & gauge, \texttt{state} & promotions by state, plus
    \texttt{refused} \\
  \texttt{pii\_memory\_entities} & gauge & entities in the memory graph \\
  \texttt{pii\_memory\_relationships} & gauge & relationships in the memory graph \\
  \texttt{pii\_memory\_writes\_total} & counter, \texttt{kind} & writes that changed the graph \\
  \texttt{pii\_memory\_write\_refusals\_total} & counter, \texttt{reason} & memory writes refused \\
  \texttt{pii\_memory\_query\_refusals\_total} & counter, \texttt{reason} & memory queries refused \\
  \texttt{pii\_memory\_tier} & gauge, \texttt{tier} & records in each tier \\
  \texttt{pii\_memory\_tier\_changes\_total} & counter, \texttt{change} & transitions,
    \texttt{<from>\_to\_<to>} \\
  \texttt{pii\_memory\_archive\_ops\_total} & counter, \texttt{op} & archived, restored, refused \\
  \bottomrule
\end{longtable}
}

Two of these are instruments rather than descriptions. \texttt{pii\_store\_bytes\_per\_run} is
the growth rate as one number: a store that is not growing prints a small value, and a value
that climbs after every run names the subsystem to look at.
\texttt{pii\_quarantine\_refusals\_recorded\_total} exists to be compared with
\texttt{pii\_quarantine\_refusals\_total}: a refusal that did not reach the ledger is a defect
in the refusal path, not in the run.

\section{Health}

\texttt{pii health} prints the identity-bearing parts of the status payload, in a fixed order:

\begin{lstlisting}[language=bash]
# the shape of the output; values abbreviated
pii 0.2.0 (c05f3d1e4dd3, x86_64-unknown-linux-gnu)
client: c05f3d1e4dd3
config: /root/.pii/pii.toml sha256=9f2c...
stores:
  main: /run/user/0/pii/store.redb (1048576 bytes, format_version=1)
mounts:
  workspace [rw]: /projects/pi-improved  (workspace)
since_deploy:
  revision: c05f3d1e4dd3
  runs: 42
  starts: 3
  hard_deaths: 1
  runs_by_status.succeeded: 39
counts:
  runs: 42
\end{lstlisting}

Read it in three passes. \emph{Identity:} the binary's revision and this client's must match;
\texttt{config} names the file that is actually in force, with its digest. \emph{State:} the
stores and their format versions, then the mounts, which are the sandbox's answer to
\enquote{what may this daemon write}. \emph{History:} \texttt{since\_deploy} is the run of
numbers that survives a restart --- \texttt{starts} against \texttt{hard\_deaths} is the
crash record, and a \texttt{hard\_deaths} that grew since yesterday is the first thing to
explain.

\begin{center}
\begin{minipage}{\linewidth}
\centering
\resizebox{\linewidth}{!}{%
\begin{tikzpicture}[node distance=8mm and 12mm]
  \node[b] (h) {\texttt{pii health}};
  \node[v, right=of h] (rev) {same revision?};
  \node[b, above right=6mm and 12mm of rev] (ok) {exit 0};
  \node[b, right=of rev, yshift=9mm] (drift) {exit 3\\\scriptsize DRIFT: not this tree};
  \node[b, right=of rev, yshift=-9mm] (unk) {exit 3\\\scriptsize UNVERIFIED: unknown revision};
  \node[b, below=of rev] (down) {exit 2\\\scriptsize no daemon answered};
  \draw[f] (h) -- (rev);
  \draw[f] (rev) -- node[p, above=2pt, midway]{yes} (ok);
  \draw[f] (rev) -- node[p, above=2pt, midway]{no} (drift);
  \draw[f] (rev) -- node[p, below=2pt, midway]{unknown} (unk);
  \draw[bk] (h) -- (down);
\end{tikzpicture}}
\captionof{figure}{Three outcomes, two of them exit 3 for different reasons. \enquote{Unknown} is never treated as agreement: a build that cannot prove what it is has not proven it is current.}
\end{minipage}
\end{center}

The distinction between drift and unknown matters in an upgrade. Drift means the daemon is a
different, known binary: restart it with the new one. Unknown means one side has no revision
stamped, which is what a locally built binary without the build script looks like --- the
comparison cannot be made, and \texttt{pii health} will not pretend otherwise.

\section{The run journal}

The journal is the only account of what a run did, and it is append-only: every event is
written once, under \texttt{(run, seq)}. Seventeen event variants exist; the ones an operator
reads are these.

\begin{center}
\begin{tabular}{@{}>{\raggedright\arraybackslash}p{0.24\textwidth}>{\raggedright\arraybackslash}p{0.64\textwidth}@{}}
  \toprule
  \textbf{Journal event} & \textbf{What it records} \\
  \midrule
  \texttt{Started} & the graph name and a hash of the run input \\
  \texttt{TickStarted} & the batch of ready nodes drained in one scheduling tick, and its width \\
  \texttt{NodeStarted} & a node began, with its millisecond offset \\
  \texttt{NodeFinished} & the port it took, the value it produced, the value's hash, the
    duration \\
  \texttt{NodeFailed} & the error, whether it is retryable, and the attempt number \\
  \texttt{NodeProgress} & a heartbeat: elapsed time and child CPU. A wedged node shows elapsed
    growing while CPU stays flat. \\
  \texttt{NodeCancelled} & a node stopped because a person cancelled the run \\
  \texttt{NodeDead} & a node marked dead; the run cannot pass it \\
  \texttt{NodeInterrupted} & the node was cut short by a suspension or shutdown \\
  \texttt{Suspended}, \texttt{Resumed} & the run parked on a resume key, and continued \\
  \texttt{NodeRetried} & an explicit retry of one node \\
  \texttt{EffectStarted}, \texttt{EffectConfirmed} & a side effect began, and was confirmed \\
  \texttt{RunSpawned} & a child run was started, with its id \\
  \texttt{Finished} & the terminal outcome \\
  \bottomrule
\end{tabular}
\end{center}

\begin{center}
\begin{minipage}{\linewidth}
\centering
\resizebox{\linewidth}{!}{%
\begin{tikzpicture}[node distance=12mm and 16mm]
  \node[b] (run) at (0,0) {running};
  \node[b] (sus) at (3.6,0) {suspended};
  \node[b] (done) at (7.0,0) {succeeded};
  \node[b] (ref) at (0,-2.6) {refused};
  \node[b] (fail) at (3.6,-2.6) {failed};
  \draw[f] (run) -- node[p, above=2pt, midway]{gate, timer} (sus);
  \draw[f] (sus) -- node[p, above=2pt, midway]{resume: pass} (done);
  \draw[bk] (sus.south) to[out=250,in=290] node[p, below=6pt, midway]{resume} (run.south);
  \draw[f] (sus) -- node[p, right=2pt, midway]{fail} (fail);
  \draw[sf] (run) -- node[p, left=6pt, midway]{pre-spawn} (ref);
\end{tikzpicture}}
\captionof{figure}{Five statuses, and the one that is easy to miss. \texttt{Refused} means the agent call was refused before anything was spawned: it is not a failure of the work, and it is not in the \texttt{run ls -{}-status} vocabulary.}
\end{minipage}
\end{center}

That last point is a real limitation of the filter. \texttt{run ls -{}-status} accepts
\texttt{running}, \texttt{suspended}, \texttt{succeeded} and \texttt{failed}; the store's
fifth status, \texttt{refused}, cannot be named on the command line. A refused run therefore
appears in an unfiltered listing and disappears from every filtered one, which is exactly
when someone is looking for it.

\begin{lstlisting}[language=bash]
pii run ls                       # refused runs appear only here
pii run trace r000000000261 --reason implement
pii run trace r000000000261 --overlapped code,docs
pii run trace r000000000261 --terminal
\end{lstlisting}

\texttt{-{}-overlapped} is the query that answers a performance question without a profiler:
two nodes that were supposed to be sequential and were not show up as overlapped, and two
that were supposed to run together and never did show up as disjoint.

\section{The audit trail}

\texttt{audit.log} sits beside the store and receives only the security-relevant events. It is
rotated like the log, it is on by default, and it exists so that an operator can answer
\enquote{what was denied} without having enabled a full event log.

\begin{center}\footnotesize
\begin{tabular}{@{}>{\raggedright\arraybackslash}p{0.34\textwidth}>{\raggedright\arraybackslash}p{0.54\textwidth}@{}}
  \toprule
  \textbf{Kind in the trail} & \textbf{What it records} \\
  \midrule
  \texttt{core.cap.denied} & a capability the node did not hold \\
  \texttt{tools.deny\_list.hit} & a tool call matched the deny list \\
  \texttt{fetch.allowlist.denied} & a fetch to a host outside the allowlist \\
  \texttt{net.egress.denied} & egress outside \texttt{[net] grants} \\
  \texttt{daemon.sandbox.in\_force} & the profile, the exact grant, and whether confinement is
    applied \\
  \texttt{daemon.peer.refused} & a control peer refused at the socket \\
  \texttt{daemon.observer.detached} & the observer fan-out was switched off \\
  \texttt{daemon.audit.disabled} & the audit feed was switched off \\
  \bottomrule
\end{tabular}
\end{center}

A start writes \texttt{daemon.sandbox.in\_force} to the trail. It carries \texttt{profile},
\texttt{grant} and \texttt{enforced}, and it is the only place that states what the daemon was
actually confined to rather than what it was asked for.

One limit is structural and stated in the module that implements the filter: a denial raised
by the kernel inside a child process --- a raw redirect blocked by Landlock --- never passes
through pii and is therefore not in the trail. Only denials pii itself decides appear.

\begin{lstlisting}[language=bash]
grep daemon.sandbox.in_force /var/lib/pii/audit.log | tail -1
grep -c core.cap.denied /var/lib/pii/audit.log
grep net.egress.denied /var/lib/pii/audit.log | jq -r '.fields.host' | sort | uniq -c
\end{lstlisting}

With \texttt{-{}-audit off} the file is not empty: it records exactly one kind,
\texttt{daemon.audit.disabled}, so a trail that stops is distinguishable from a trail that was
never written.

\section{Incidents and the loop}

Two surfaces cover the unattended side, and neither needs a model.

\texttt{pii incidents scan -{}-source runtime -{}-log F} reads the daemon's own log and reports
each incident with its count, first and last seen, and \texttt{latency\_secs} against the
configured window --- the detection latency, which is the number that says whether the alert
arrived in time. Adding \texttt{-{}-ingest -{}-dir D} stores what it found;
\texttt{pii incidents show} reads one back. \texttt{pii incidents distance -{}-since} measures
the operator's own response time from recorded timestamps.

\texttt{pii loop emit} writes one NDJSON line per lifecycle step of an autonomy loop, and
\texttt{pii loop events} merges that file with the daemon log so the loop and the daemon read
as one timeline. The exit code carries information the text does not: \texttt{emit} keeps the
event locally and exits 3 when the daemon did not answer, appending \texttt{daemon=down} to
the detail. A loop that records 3s is a loop that was running while the daemon was not.

\section{Correlating a failure}

The order matters, because each step narrows the next.

\begin{center}
\begin{minipage}{\linewidth}
\centering
\resizebox{\linewidth}{!}{%
\begin{tikzpicture}[node distance=8mm and 12mm]
  \node[b] (one) {1. \texttt{pii health}};
  \node[b, right=of one] (two) {2. \texttt{pii metrics}};
  \node[b, right=of two] (three) {3. \texttt{run ls}};
  \node[b, below=of three] (four) {4. \texttt{run trace}};
  \node[b, left=of four] (five) {5. the log, by \texttt{run}};
  \node[b, left=of five] (six) {6. \texttt{audit.log}};
  \draw[f] (one) -- (two);
  \draw[f] (two) -- (three);
  \draw[f] (three) -- (four);
  \draw[f] (four) -- (five);
  \draw[f] (five) -- (six);
  \draw[bk] (six.south) to[bend left=14] node[p, below=2pt, midway]{a denial explains it} (one.south);
\end{tikzpicture}}
\captionof{figure}{Six steps from \enquote{something is wrong} to a cause. Steps 1 and 2 rule out the daemon; steps 3 to 5 locate the run and the node; step 6 explains a refusal that the run itself could not.}
\end{minipage}
\end{center}

\begin{longtable}{@{}>{\raggedright\arraybackslash}p{0.30\textwidth}>{\raggedright\arraybackslash}p{0.58\textwidth}@{}}
  \toprule
  \textbf{Symptom} & \textbf{First surface, and what to look for} \\
  \midrule
  \endhead
  every run fails in a second & the log, for \texttt{model.provider.auth\_failed} or
    \texttt{rate\_limited}; the metric set carries no model series, so this is an event
    question, not a metric one \\
  a run is \enquote{running} for hours & \texttt{run trace -{}-node-status}, then
    \texttt{NodeProgress}: elapsed growing with flat CPU is a wedged child \\
  the run stopped with no error & \texttt{run trace -{}-terminal}: a node left through a port
    with no edge, so nothing failed and nothing continued \\
  a write was refused & \texttt{audit.log}: \texttt{core.cap.denied} or
    \texttt{tools.deny\_list.hit}, never the run's own output \\
  the log has a gap & \texttt{observe.events.dropped}: the pipeline dropped events, so absence
    is not evidence \\
  the daemon restarted by itself & \texttt{pii health}: \texttt{hard\_deaths} in
    \texttt{since\_deploy}, and a log that ends on \texttt{daemon.started} with no
    \texttt{daemon.stopped} \\
  the metrics do not move & the daemon serving them is not the one running the work: compare
    \texttt{pii\_uptime\_ms} with the process you restarted \\
  \bottomrule
\end{longtable}

\begin{lstlisting}[language=bash]
pii health --json > /tmp/health.json
pii metrics | grep -E 'provider|store_bytes|latency'
pii run ls --status failed | jq -r '.[].run' | head
pii run trace "$RUN" --reason "$NODE"
grep "\"run\":\"$RUN\"" /var/log/pii/pii.log | jq -c 'select(.level!="info")'
grep -E 'cap.denied|egress.denied' /var/lib/pii/audit.log | tail -20
\end{lstlisting}